Since the release of ChatGPT in November 2022, large language models (LLMs) have been adopted widely across software engineering (SE) research~\cite{DBLP:journals/tosem/HouZLYWLLLGW24}, yet the reproducibility and replicability of empirical studies involving LLMs remains uncertain.
% GROUNDING DBLP:journals/tosem/HouZLYWLLLGW24: "Large Language Models (LLMs) have significantly impacted numerous domains, including Software Engineering (SE). Many recent publications have explored LLMs applied to various SE tasks."
Recent findings indicate a low reproducibility of SE studies involving LLMs~\cite{DBLP:journals/corr/abs-2510-25506}. \citeauthor{DBLP:conf/ndss/EvertzRNMNSGPSS26} surveyed 72 peer-reviewed papers from leading security and SE venues published in 2023--2024 and found that every paper contained at least one of nine LLM-specific pitfalls they identified, spanning data collection, pre-training, fine-tuning, prompting, and evaluation; only 15.7\% of the pitfalls present in those papers were explicitly discussed~\cite{DBLP:conf/ndss/EvertzRNMNSGPSS26}.
% GROUNDING DBLP:journals/corr/abs-2510-25506: "For none of the five studies, we were able to fully reproduce the results."
% GROUNDING DBLP:conf/ndss/EvertzRNMNSGPSS26: "all 72 peer-reviewed papers published at leading Security and Software Engineering venues between 2023 and 2024, finding that every paper contains at least one pitfall"
Three characteristics of LLMs make this particularly challenging.
First, their inherent non-determinism causes variability across runs~\cite{DBLP:conf/naacl/SongWLL25, DBLP:journals/corr/abs-2602-07150}, and slight changes can lead to substantially different results~\cite{DBLP:conf/nips/AgarwalSCCB21, DBLP:journals/corr/abs-2412-12509}.
% GROUNDING DBLP:conf/naacl/SongWLL25: "given a particular input, the same LLM may generate significantly different outputs under various decoding configurations, a phenomenon known as non-determinism in generation"
% GROUNDING DBLP:journals/corr/abs-2602-07150: "single-run pass@1 estimates vary by 2.2 to 6.0 percentage points depending on which run is selected, with standard deviations exceeding 1.5 percentage points even at temperature 0"
% GROUNDING DBLP:conf/nips/AgarwalSCCB21: "A small number of training runs coupled with high variability in performance of deep RL algorithms often leads to substantial statistical uncertainty in reported point estimates."
% GROUNDING DBLP:journals/corr/abs-2412-12509: "A 2% variation in the LLM's judgment could lead to significant shifts in the ranking, potentially misrepresenting the true performance of the models"
Second, commercial models evolve beyond version identifiers, so reported performance can change over time~\cite{DBLP:journals/corr/abs-2307-09009}.
% GROUNDING DBLP:journals/corr/abs-2307-09009: "We find that the performance and behavior of both GPT-3.5 and GPT-4 can vary greatly over time."
Third, even for ``open'' models, training data and fine-tuning details often remain undisclosed~\cite{Gibney2024}.
% GROUNDING Gibney2024: "Many of the large language models that power chatbots claim to be open, but restrict access to code and training data."
Moreover, prompt formatting choices alone can shift accuracy by up to 76 percentage points~\cite{DBLP:conf/iclr/Sclar0TS24}, and configured parameters such as temperature affect output variability~\cite{renze2024temperature}.
% GROUNDING DBLP:conf/iclr/Sclar0TS24: "several widely used open-source LLMs are extremely sensitive to subtle changes in prompt formatting in few-shot settings, with performance differences of up to 76 accuracy points"
% GROUNDING renze2024temperature: "A higher temperature makes the output more random, thus favoring more 'creative' predictions."
Hence, not reporting these settings directly affects reproducibility.
Traditional open science practice in SE has focused on releasing source code and datasets as a replication package. The empirical artifact analysis literature in SE focuses more on code repositories and data than on upstream artifacts such as requirements specifications and design documents~\cite{DBLP:journals/jss/LiuHHXZJCM24}. With LLMs, code is often generated from prompts, context files, and other runtime inputs to the model, so reporting must shift left to cover these upstream artifacts in addition to the code and data they produce.
% GROUNDING DBLP:journals/jss/LiuHHXZJCM24: "All the top-starred artifacts include executable code, providing either tools, frameworks, or toolkits."

Although the SE research community has developed guidelines for conducting and reporting specific types of empirical studies such as controlled experiments (e.g.,~\cite{DBLP:books/sp/WohlinRHORW24, DBLP:books/sp/08/SSS2008}), their replications (e.g.,~\cite{DBLP:journals/tse/SantosVOJ21}), or empirical studies in general (e.g., the \emph{ACM SIGSOFT Empirical Standards}~\cite{ralph2021empiricalstandardssoftwareengineering}), none of these address the LLM-specific aspects described above.
% GROUNDING DBLP:books/sp/WohlinRHORW24: "Part II then devotes one chapter to each of the five experiment steps: scoping, planning, execution, analysis, and result presentation."
% GROUNDING DBLP:books/sp/08/SSS2008: "focus on the practical knowledge necessary for conducting, reporting, and using empirical methods in software engineering"
% GROUNDING DBLP:journals/tse/SantosVOJ21: "Provide an analysis procedure with a set of embedded guidelines to aggregate software engineering (SE) replication results."
% GROUNDING ralph2021empiricalstandardssoftwareengineering: "There are three kinds of standards. The General Standard applies to all empirical research."
Previously, a position paper highlighted these issues~\cite{DBLP:conf/wsese/0001BFB25}, but there was no comprehensive community-developed guidance for designing and reporting empirical studies involving LLMs in SE.
% GROUNDING DBLP:conf/wsese/0001BFB25: "While providing a comprehensive set of guidelines is beyond the scope of this position paper, we report a first set of guidelines based on a discussion session with other empiricism experts"

Therefore, we present community-developed guidelines for designing and reporting studies involving LLMs in SE research, co-developed by 22 researchers.
After outlining our \scope, we introduce a taxonomy of \studytypes, then present eight \guidelines.
We complement these with an applicability matrix mapping guidelines to study types and a reporting checklist for authors and reviewers.
For each study type and guideline, we identify relevant examples, both within and outside of SE research.
We maintain the study types and guidelines online as a living resource for the community to use and shape (\href{https://llm-guidelines.org/}{llm-guidelines.org}).
